\documentclass{apcalc}
\title{AP Calculus BC: Full Practice Exam}
\author{Paul Wang}
\date{November 2026}
\volume{May 2027 format}
\shorttitle{Full Practice Exam}
\begin{document}
\maketitle
\begin{abstract}
A full-length practice exam in the May 2027 format: Section I has 42 multiple-choice questions (Part A, 29 questions in 62 minutes without a calculator; Part B, 13 questions in 38 minutes with a graphing calculator), and Section II has six nine-point free-response questions (Part A, 2 questions in 30 minutes with a calculator; Part B, 4 questions in 60 minutes without). The multiple-choice questions follow the unit weighting in the course description, and every item is tagged with its unit and the rule from the unit guides that it tests. All questions are original. The key at the end gives an answers-only strip, a reason and rule for every multiple-choice item, a scoring worksheet, and the free-response scoring guides.
\end{abstract}
\scoringline{42 MC \quad\blank\ / 54 FR}{Section I: 100 minutes. Section II: 90 minutes. Take the sections in order, timed, with a 10-minute break between them.}
\section*{Directions}
Section I: choose the best of the four answers and write its letter in the blank beside the question. There is no penalty for wrong answers. In Part B, if the exact value is not one of the choices, select the choice that best approximates it. Section II: show all work. Justify every conclusion with calculus, write the integral or expression you evaluate on the calculator, and give decimal answers correct to three places. Unless otherwise specified, the domain of a function is the set of real numbers for which it is defined, and angles are in radians.
\clearpage\section*{Section I, Part A: 29 questions, 62 minutes. No calculator.}
\begin{q}{1}
$\displaystyle\lim_{x\to3}\frac{x^2-9}{x^2-2x-3}$ is
\chfour{$0$}{$1$}{$\dfrac32$}{nonexistent}
\end{q}
\begin{q}{2}
$\displaystyle\lim_{x\to\infty}\frac{3x^2-5}{7+x-6x^2}$ is
\chfour{$-\dfrac12$}{$0$}{$\dfrac12$}{$\dfrac37$}
\end{q}
\begin{q}{3}
$\displaystyle\lim_{h\to0}\frac{(2+h)^4-16}{h}$ is
\chfour{$16$}{$24$}{$32$}{$64$}
\end{q}
\begin{q}{4}
Let $f(x)=\begin{cases}x^2+a,&x\le1\\ bx,&x>1\end{cases}$. For which $a$ and $b$ is $f$ differentiable at $x=1$?
\chtwo{$a=0$, $b=2$}{$a=1$, $b=2$}{$a=2$, $b=3$}{$a=1$, $b=1$}
\end{q}
\begin{q}{5}
If $f(x)=x\sec x$, then $f'(0)=$
\chfour{$0$}{$1$}{$2$}{undefined}
\end{q}
\begin{q}{6}
Let $h(x)=g(x^2)$, where $g$ is differentiable and $g'(4)=-3$. What is $h'(2)$?
\chfour{$-12$}{$-6$}{$-3$}{$12$}
\end{q}
\begin{q}{7}
If $y=\arctan(\sqrt x)$, then $\dfrac{dy}{dx}$ at $x=1$ is
\chfour{$\dfrac14$}{$\dfrac12$}{$1$}{$\dfrac\pi4$}
\end{q}
\begin{q}{8}
The area of a circle increases at $6\pi$ cm² per second. How fast is the radius increasing when the radius is 3 cm?
\chfour{$\dfrac13$ cm/s}{$1$ cm/s}{$3$ cm/s}{$6\pi$ cm/s}
\end{q}
\begin{q}{9}
$\displaystyle\lim_{x\to0}\frac{e^{3x}-1}{\sin x}$ is
\chfour{$0$}{$1$}{$3$}{nonexistent}
\end{q}
\begin{q}{10}
If $f'(x)=x(x-2)^2(x+3)$, then $f$ has a relative maximum at
\chfour{$x=-3$}{$x=0$}{$x=2$}{$x=-3$ and $x=2$}
\end{q}
\begin{q}{11}
If $f''(x)=(x-1)(x-4)^2$, the graph of $f$ has points of inflection at
\chtwo{$x=4$ only}{$x=1$ and $x=4$}{no values of $x$}{$x=1$ only}
\end{q}
\begin{q}{12}
The graph of $f'$ is shown for $0\le x\le5$. On which interval is $f$ both decreasing and concave up?
\begin{center}\includegraphics[width=0.45\linewidth]{figs/ex-q12.pdf}\end{center}
\chfour{$(1,2.5)$}{$(1,4)$}{$(4,5)$}{$(2.5,4)$}
\end{q}
\begin{q}{13}
The left Riemann sum for $\int_0^3x^2\,dx$ with three equal subintervals is
\chfour{$5$}{$9$}{$9.5$}{$14$}
\end{q}
\begin{q}{14}
If $F(x)=\displaystyle\int_0^{2x}\sin(t^2)\,dt$, then $F'(x)=$
\chfour{$\sin(4x^2)$}{$2\sin(2x^2)$}{$4x\cos(4x^2)$}{$2\sin(4x^2)$}
\end{q}
\begin{q}{15}
$\displaystyle\int_0^{\pi/4}\tan x\sec^2x\,dx=$
\chfour{$\dfrac14$}{$\dfrac12$}{$1$}{$\dfrac{\sqrt2}{2}$}
\end{q}
\begin{q}{16}
$\displaystyle\int x\ln x\,dx=$
\chtwo{$x\ln x-x+C$}{$\dfrac{x^2}{2}\ln x+\dfrac{x^2}{4}+C$}{$\ln x+1+C$}{$\dfrac{x^2}{2}\ln x-\dfrac{x^2}{4}+C$}
\end{q}
\begin{q}{17}
$\displaystyle\int_2^3\frac{dx}{x^2-1}=$
\chfour{$\dfrac12\ln\dfrac32$}{$\ln\dfrac32$}{$\dfrac12\ln2$}{$\ln\dfrac43$}
\end{q}
\begin{q}{18}
$\displaystyle\int_1^\infty xe^{-x^2}\,dx$ is
\chfour{$\dfrac{1}{2e}$}{$\dfrac1e$}{$\dfrac12$}{divergent}
\end{q}
\begin{q}{19}
Let $y=f(x)$ solve $\dfrac{dy}{dx}=xy$ with $f(0)=2$. Euler's method with two steps of size $0.5$ gives which approximation of $f(1)$?
\chfour{$2$}{$2.25$}{$2.5$}{$3$}
\end{q}
\begin{q}{20}
The solution of $\dfrac{dy}{dx}=y\cos x$ with $y(0)=3$ is
\chfour{$y=e^{\sin x}+2$}{$y=3e^{\cos x}$}{$y=3+\sin x$}{$y=3e^{\sin x}$}
\end{q}
\begin{q}{21}
The area of the region enclosed by $y=4-x^2$ and $y=x+2$ is
\chfour{$\dfrac76$}{$\dfrac32$}{$\dfrac{10}{3}$}{$\dfrac92$}
\end{q}
\begin{q}{22}
For $x=e^t$, $y=t^2+1$, the value of $\dfrac{dy}{dx}$ at $t=1$ is
\chfour{$\dfrac2e$}{$\dfrac e2$}{$2$}{$2e$}
\end{q}
\begin{q}{23}
A particle has velocity $\langle6t,\,3t^2\rangle$. How far does it travel for $0\le t\le1$?
\chfour{$\sqrt{10}$}{$5\sqrt5-8$}{$4$}{$5\sqrt5$}
\end{q}
\begin{q}{24}
The area of the region enclosed by the polar curve $r=3+2\cos\theta$ is
\chfour{$\dfrac{11\pi}{2}$}{$11\pi$}{$13\pi$}{$22\pi$}
\end{q}
\begin{q}{25}
$\displaystyle\sum_{n=1}^{\infty}\frac{4^n}{5^{n+1}}=$
\chfour{$\dfrac{4}{45}$}{$\dfrac{4}{25}$}{$\dfrac45$}{$4$}
\end{q}
\begin{q}{26}
Which series converges?
\chfour{$\displaystyle\sum\frac{n!}{5^n}$}{$\displaystyle\sum\frac{n+1}{n}$}{$\displaystyle\sum\frac{1}{n^{1/3}}$}{$\displaystyle\sum\frac{5^n}{n!}$}
\end{q}
\begin{q}{27}
The interval of convergence of $\displaystyle\sum_{n=1}^\infty\frac{(x-3)^n}{n\,5^n}$ is
\chfour{$(-2,8)$}{$(-2,8]$}{$[-2,8)$}{$[-2,8]$}
\end{q}
\begin{q}{28}
$f(0)=1$, $f'(0)=2$, $f''(0)=-6$, $f'''(0)=12$. The third-degree Taylor polynomial for $f$ about $x=0$ is
\chtwo{$1+2x-6x^2+12x^3$}{$1+2x-3x^2+4x^3$}{$1+2x-6x^2+2x^3$}{$1+2x-3x^2+2x^3$}
\end{q}
\begin{q}{29}
In the Maclaurin series for $\cos(2x)$, the coefficient of $x^4$ is
\chfour{$-\dfrac23$}{$\dfrac{1}{24}$}{$\dfrac23$}{$\dfrac43$}
\end{q}
\clearpage\section*{Section I, Part B: 13 questions, 38 minutes. Graphing calculator required.}
\begin{q}{30}
How many solutions does $x^5-3x+1=0$ have on $[-2,2]$?
\chfour{One}{Two}{Three}{Five}
\end{q}
\begin{q}{31}
Let $f(x)=x^5+x+1$ and $g=f^{-1}$. What is $g'(5)$?
\chfour{$0.0003$}{$0.081$}{$1.226$}{$12.309$}
\end{q}
\begin{q}{32}
A particle moves on the $x$-axis with velocity $v(t)=t\sin(t^2)$. Which is true at $t=2$?
\chone{$a(2)\approx-5.986$, and the speed is increasing}{$a(2)\approx-5.986$, and the speed is decreasing}{$a(2)\approx-1.514$, and the speed is increasing}{$a(2)\approx5.986$, and the speed is decreasing}
\end{q}
\begin{q}{33}
What is the absolute maximum value of $f(x)=\sin x+0.3x^2-x$ on $[1,4]$?
\chfour{$0.043$}{$0.142$}{$0.174$}{$1.426$}
\end{q}
\begin{q}{34}
The derivative of $f$ is $f'(x)=e^{\cos x}-\frac x2$ for $0<x<6$. At what value of $x$ does $f$ have a relative maximum?
\chfour{$1.721$}{$3.142$}{$4.155$}{$5.000$}
\end{q}
\begin{q}{35}
$F(0)=10$ and $F'(t)=\ln(1+t^2)$. What is $F(3)$?
\chfour{$3.406$}{$12.303$}{$13.406$}{$16.908$}
\end{q}
\begin{q}{36}
For what value of $k$ is $\displaystyle\int_0^ke^{-x^2}\,dx=0.8$?
\chfour{$0.472$}{$0.800$}{$1.172$}{$1.609$}
\end{q}
\begin{q}{37}
A population satisfies $\dfrac{dP}{dt}=0.02P\left(1-\dfrac{P}{800}\right)$ with $P(0)=100$, so $P(t)=\dfrac{800}{1+7e^{-0.02t}}$. What is $P(50)$?
\chfour{$223.766$}{$271.828$}{$400.000$}{$800.000$}
\end{q}
\begin{q}{38}
What is the average value of $f(x)=\sqrt{1+x^3}$ on $[0,2]$?
\chfour{$1.000$}{$1.621$}{$2.000$}{$3.241$}
\end{q}
\begin{q}{39}
The region between $y=2\sin x$ and $y=x$ in the first quadrant is revolved about the $x$-axis. What is the volume?
\chfour{$0.842$}{$1.433$}{$2.126$}{$6.678$}
\end{q}
\begin{q}{40}
A particle has $x'(t)=e^{-t}+\sin(t^2)$ and $x(0)=2$. What is $x(2)$?
\chfour{$1.379$}{$1.669$}{$2.865$}{$3.669$}
\end{q}
\begin{q}{41}
What is the area of the region bounded by $r=2+\sin(2\theta)$ for $0\le\theta\le\frac\pi2$ and the coordinate axes?
\chfour{$2.071$}{$5.534$}{$7.069$}{$11.069$}
\end{q}
\begin{q}{42}
The second-degree Taylor polynomial for $f$ about $x=1$ is $3-2(x-1)+5(x-1)^2$, and $|f'''(x)|\le12$ for $1\le x\le1.3$. The Lagrange error bound for $|f(1.3)-P_2(1.3)|$ is
\chfour{$0.027$}{$0.054$}{$0.162$}{$0.324$}
\end{q}
\clearpage\section*{Section II, Part A: 2 questions, 30 minutes. Graphing calculator required.}
Show all work. Justify your answers, and give decimal answers correct to three places.
\begin{frq}{Question 1 --- Water in a tank (9 points, calculator required)}
A tank holds 150 gallons of water at time $t=0$. For $0\le t\le12$ hours, water is pumped into the tank at the rate $W(t)=25e^{-0.05t}$ gallons per hour and drains out at the rate $D(t)=20+10\sin\left(\frac t2\right)$ gallons per hour.
\partlabel{a}How many gallons of water are pumped into the tank during $0\le t\le12$?\workspace{1.6in}
\partlabel{b}Is the amount of water in the tank increasing or decreasing at $t=6$? Give a reason for your answer.\workspace{1.6in}
\partlabel{c}At what time $t$, $0\le t\le12$, is the amount of water in the tank least? Justify your answer.\workspace{1.6in}
\partlabel{d}Find $D'(4)$. Using correct units, explain the meaning of this value in context.\workspace{1.6in}
\end{frq}
\begin{frq}{Question 2 --- A polar curve (9 points, calculator required)}
The curve $r=\theta+\sin(2\theta)$ for $0\le\theta\le\pi$ is drawn in the $xy$-plane.
\partlabel{a}Find the area of the region bounded by the curve and the $x$-axis.\workspace{1.6in}
\partlabel{b}Find $\frac{dr}{d\theta}$ at $\theta=\frac\pi2$. Is the point on the curve moving toward or away from the origin as $\theta$ increases through $\frac\pi2$? Explain.\workspace{1.6in}
\partlabel{c}Find the slope of the curve at $\theta=\frac\pi2$.\workspace{1.6in}
\partlabel{d}Find the greatest distance from the origin to a point on the curve. Justify your answer.\workspace{1.6in}
\end{frq}
\clearpage\section*{Section II, Part B: 4 questions, 60 minutes. No calculator.}
\begin{frq}{Question 3 --- A car's velocity (9 points, no calculator)}
A car moves along a straight road. Its velocity $v(t)$, in meters per second, is differentiable; selected values are given in the table.
\begin{center}\small\renewcommand{\arraystretch}{1.35}\begin{tabular}{|c|c|c|c|c|c|}
\hline
$t$ (seconds) & 0 & 2 & 5 & 8 & 12 \\ \hline
$v(t)$ (meters per second) & 10 & 16 & 22 & 18 & 6 \\ \hline
\end{tabular}\end{center}
\partlabel{a}Use a trapezoidal sum with the four subintervals in the table to approximate $\int_0^{12}v(t)\,dt$. Explain the meaning of this integral in context.\workspace{1.6in}
\partlabel{b}Must there be a time $t$, $2<t<5$, at which $v(t)=20$? Justify your answer.\workspace{1.6in}
\partlabel{c}Must there be a time $t$, $8<t<12$, at which the acceleration of the car is $-3$ meters per second per second? Justify your answer.\workspace{1.6in}
\partlabel{d}For $12\le t\le15$ the car slows with constant acceleration $-2$ meters per second per second. Find $v(15)$ and the distance the car travels from $t=12$ to $t=15$.\workspace{1.6in}
\partlabel{e}Using your answer from part (a), estimate the average velocity of the car over $0\le t\le12$.\workspace{1.6in}
\end{frq}
\begin{frq}{Question 4 --- An accumulation function (9 points, no calculator)}
The continuous function $f$ is defined on $[-3,4]$; its graph consists of four line segments, as shown. Let $g(x)=5+\int_0^xf(t)\,dt$.
\begin{center}\includegraphics[width=0.55\linewidth]{figs/ex-fr4.pdf}\end{center}
\partlabel{a}Find $g(-3)$ and $g(4)$.\workspace{1.6in}
\partlabel{b}Find the $x$-coordinate of each relative extremum of $g$ on $(-3,4)$ and classify it. Justify your answers.\workspace{1.6in}
\partlabel{c}Find the absolute minimum value of $g$ on $[-3,4]$. Justify your answer.\workspace{1.6in}
\partlabel{d}Find $\displaystyle\lim_{x\to3}\frac{g(x)-10}{(x-3)^2}$.\workspace{1.6in}
\end{frq}
\begin{frq}{Question 5 --- A differential equation (9 points, no calculator)}
Let $y=f(x)$ be the particular solution to $\dfrac{dy}{dx}=y^2(2x+1)$ with $f(0)=-1$.
\partlabel{a}Write an equation for the line tangent to the graph of $f$ at $x=0$, and use it to approximate $f(0.2)$.\workspace{1.6in}
\partlabel{b}Use Euler's method, starting at $x=0$ with two steps of equal size, to approximate $f(0.2)$.\workspace{1.6in}
\partlabel{c}Find $y=f(x)$, the particular solution with $f(0)=-1$.\workspace{1.6in}
\end{frq}
\begin{frq}{Question 6 --- A Taylor polynomial (9 points, no calculator)}
A function $f$ has derivatives of all orders, with $f(0)=1$, $f'(0)=-2$, $f''(0)=6$, $f'''(0)=-24$, and $|f^{(4)}(x)|\le60$ for $0\le x\le0.5$.
\partlabel{a}Write the third-degree Taylor polynomial $P_3(x)$ for $f$ about $x=0$.\workspace{1.6in}
\partlabel{b}Use $P_3$ to approximate $f(0.5)$, and use the Lagrange error bound to show that $|f(0.5)-P_3(0.5)|<0.16$.\workspace{1.6in}
\partlabel{c}Let $g(x)=\int_0^xf(t)\,dt$. Write the fourth-degree Taylor polynomial for $g$ about $x=0$.\workspace{1.6in}
\partlabel{d}Let $h(x)=f(x^2)$. Write the first three nonzero terms of the Taylor series for $h$ about $x=0$.\workspace{1.6in}
\end{frq}
\clearpage\section*{Answer key}
\keystrip{\keyitem{1}{C}\keyitem{2}{A}\keyitem{3}{C}\keyitem{4}{B}\keyitem{5}{B}\keyitem{6}{A}\keyitem{7}{A}\keyitem{8}{B}\keyitem{9}{C}\keyitem{10}{A}\keyitem{11}{D}\keyitem{12}{D}\keyitem{13}{A}\keyitem{14}{D}\keyitem{15}{B}\keyitem{16}{D}\keyitem{17}{A}\keyitem{18}{A}\keyitem{19}{C}\keyitem{20}{D}\keyitem{21}{D}\keyitem{22}{A}\keyitem{23}{B}\keyitem{24}{B}\keyitem{25}{C}\keyitem{26}{D}\keyitem{27}{C}\keyitem{28}{D}\keyitem{29}{C}\keyitem{30}{C}\keyitem{31}{B}\keyitem{32}{A}\keyitem{33}{C}\keyitem{34}{A}\keyitem{35}{C}\keyitem{36}{C}\keyitem{37}{A}\keyitem{38}{B}\keyitem{39}{D}\keyitem{40}{D}\keyitem{41}{B}\keyitem{42}{B}}
\vspace{4pt}
\begin{keytable}
\keysection{Part A (no calculator)}
1 & C & $\frac{(x-3)(x+3)}{(x-3)(x+1)}=\frac{x+3}{x+1}\to\frac64$. & R1.5\\
2 & A & Equal degrees: $\frac{3}{-6}$. & R1.12\\
3 & C & It is $f'(2)$ for $f(x)=x^4$: $4\cdot2^3$. & R2.3\\
4 & B & Slopes: $b=2$; values: $1+a=b$. & R2.6\\
5 & B & $\sec x+x\sec x\tan x$ at 0 is $1+0$. & R2.10\\
6 & A & $h'(x)=g'(x^2)\cdot2x$; $h'(2)=g'(4)\cdot4$. & R3.2\\
7 & A & $\frac{1}{1+x}\cdot\frac{1}{2\sqrt x}=\frac12\cdot\frac12$. & R3.6\\
8 & B & $\frac{dA}{dt}=2\pi r\frac{dr}{dt}$: $6\pi=6\pi\frac{dr}{dt}$. & R4.4\\
9 & C & L'Hôpital: $\frac{3e^{3x}}{\cos x}\to3$. & R4.9\\
10 & A & $f'$ changes $+\to-$ only at $-3$. & R5.5\\
11 & D & $f''$ changes sign at 1 only; the squared factor does not change sign. & R5.7\\
12 & D & $f'<0$ on $(1,4)$; $f'$ increasing on $(2.5,5)$. & R5.9\\
13 & A & $1\cdot(0+1+4)=5$. & R6.2\\
14 & D & $\sin((2x)^2)\cdot2$. & R6.4\\
15 & B & $u=\tan x$: $\int_0^1u\,du=\frac12$. & R6.9\\
16 & D & Parts: $u=\ln x$, $dv=x\,dx$. & R6.11\\
17 & A & $\frac{1}{x^2-1}=\frac12\left(\frac{1}{x-1}-\frac{1}{x+1}\right)$. & R6.12\\
18 & A & $\lim_{b\to\infty}\left[-\frac12e^{-x^2}\right]_1^b=\frac{1}{2e}$. & R6.13\\
19 & C & Slope 0 at $(0,2)$; slope $1$ at $(0.5,2)$: $2+0.5$. & R7.4\\
20 & D & $\ln|y|=\sin x+C$; $y(0)=3$ gives $y=3e^{\sin x}$. & R7.6\\
21 & D & Intersections $x=-2,1$; $\int_{-2}^1(2-x-x^2)\,dx$. & R8.4\\
22 & A & $\frac{2t}{e^t}$ at 1. & R9.1\\
23 & B & Speed $3t\sqrt{4+t^2}$; $\int_0^1=[(4+t^2)^{3/2}]_0^1$. & R9.5\\
24 & B & $\frac12\int_0^{2\pi}(9+12\cos\theta+4\cos^2\theta)\,d\theta=\frac12(18\pi+4\pi)$. & R9.9\\
25 & C & First term $\frac{4}{25}$, ratio $\frac45$: $\frac{4/25}{1/5}$. & R10.2\\
26 & D & Ratio $\frac{5}{n+1}\to0<1$. & R10.8\\
27 & C & Radius 5 about 3; $x=-2$: AST converges; $x=8$: harmonic diverges. & R10.13\\
28 & D & Coefficients $\frac{f^{(k)}(0)}{k!}$: $-\frac62$, $\frac{12}{6}$. & R10.12\\
29 & C & $\frac{(2x)^4}{4!}=\frac{16x^4}{24}$. & R10.14\\
\keysection{Part B (calculator)}
30 & C & Graph: zeros near $-1.389$, $0.335$, $1.215$. & R1.16\\
31 & B & $f(a)=5$ at $a\approx1.2263$; $g'(5)=\frac{1}{5a^4+1}$. & R3.5\\
32 & A & $v(2)\approx-1.514<0$ and $a(2)\approx-5.986<0$: same sign. & R4.2\\
33 & C & Critical point $x\approx1.426$ with $f\approx0.174$; endpoints give $0.142$ and $0.043$. & R5.6\\
34 & A & $f'$ changes from positive to negative at $x\approx1.721$ (its only zero in $(0,6)$). & R5.5\\
35 & C & $10+\int_0^3\ln(1+t^2)\,dt\approx10+3.406$. & R6.7\\
36 & C & Solve $\int_0^ke^{-x^2}dx=0.8$ numerically. & R6.4\\
37 & A & $\frac{800}{1+7e^{-1}}\approx223.766$. & R7.8\\
38 & B & $\frac12\int_0^2\sqrt{1+x^3}\,dx\approx1.621$. & R8.1\\
39 & D & $\pi\int_0^{b}(4\sin^2x-x^2)\,dx$, $b\approx1.895$. & R8.9\\
40 & D & $2+\int_0^2(e^{-t}+\sin t^2)\,dt\approx2+1.669$. & R9.5\\
41 & B & $\frac12\int_0^{\pi/2}(2+\sin2\theta)^2\,d\theta\approx5.534$. & R9.9\\
42 & B & $\frac{12}{3!}(0.3)^3=2(0.027)$. & R10.11\\
\end{keytable}
\section*{Scoring worksheet}
Multiple choice: number correct $\times\,1.2857=$ \blank\ (out of 54). Free response: total points \blank\ (out of 54). Composite: the sum, out of 108. As a rough guide from past exams (cut scores change every year): 68+ is about a 5, 52+ a 4, 40+ a 3, 29+ a 2.
\section*{Free-response scoring guide}
\subsection*{Question 1: Water in a tank (9 points)}
\textbf{(a)} $\int_0^{12}W(t)\,dt\approx225.594$ gallons.
\begin{itemize}[leftmargin=1.6em,itemsep=0pt]
\item 1 point: integral $\int_0^{12}W(t)\,dt$
\item 1 point: answer $225.594$ gallons
\end{itemize}
\textbf{(b)} $W(6)-D(6)\approx-2.891<0$, so the amount of water is decreasing at $t=6$.
\begin{itemize}[leftmargin=1.6em,itemsep=0pt]
\item 1 point: considers $W(6)-D(6)$
\item 1 point: decreasing, because $W(6)-D(6)\approx-2.891<0$
\end{itemize}
\textbf{(c)} $A'(t)=W(t)-D(t)=0$ at $t\approx0.822$, $6.712$, $11.327$. A minimum can occur only at $t=0$, $6.712$ (where $A'$ changes from negative to positive) or $t=12$. $A(t)=150+\int_0^t(W-D)\,dx$: $A(0)=150$, $A(6.712)\approx118.763$, $A(12)\approx134.798$. The least amount is at $t\approx6.712$ hours.
\begin{itemize}[leftmargin=1.6em,itemsep=0pt]
\item 1 point: sets $W(t)-D(t)=0$ and identifies the candidates (endpoints and $t\approx6.712$)
\item 1 point: amounts at the candidates: $150$, $\approx118.763$, $\approx134.798$
\item 1 point: answer $t\approx6.712$, with the comparison as justification
\end{itemize}
\textbf{(d)} $D'(4)=5\cos2\approx-2.081$. At $t=4$ hours, the rate at which water drains is decreasing at 2.081 gallons per hour per hour.
\begin{itemize}[leftmargin=1.6em,itemsep=0pt]
\item 1 point: $D'(4)\approx-2.081$
\item 1 point: interpretation: at $t=4$ the drain rate is decreasing at 2.081 gallons per hour per hour
\end{itemize}
\subsection*{Question 2: A polar curve (9 points)}
\textbf{(a)} $\frac12\int_0^\pi(\theta+\sin2\theta)^2\,d\theta\approx4.382$.
\begin{itemize}[leftmargin=1.6em,itemsep=0pt]
\item 1 point: integral $\frac12\int_0^{\pi}r^2\,d\theta$
\item 1 point: answer $4.382$
\end{itemize}
\textbf{(b)} $\frac{dr}{d\theta}=1+2\cos2\theta=1+2\cos\pi=-1$. $r(\frac\pi2)=\frac\pi2>0$ and $r$ is decreasing, so the point is moving toward the origin.
\begin{itemize}[leftmargin=1.6em,itemsep=0pt]
\item 1 point: $\frac{dr}{d\theta}=-1$ at $\frac\pi2$
\item 1 point: toward the origin, because $r>0$ and $\frac{dr}{d\theta}<0$
\end{itemize}
\textbf{(c)} $\frac{dy}{d\theta}=r'\sin\theta+r\cos\theta=-1$ and $\frac{dx}{d\theta}=r'\cos\theta-r\sin\theta=-\frac\pi2$ at $\frac\pi2$. $\frac{dy}{dx}=\frac{-1}{-\pi/2}=\frac2\pi\approx0.637$.
\begin{itemize}[leftmargin=1.6em,itemsep=0pt]
\item 1 point: $\frac{dy}{d\theta}$ and $\frac{dx}{d\theta}$ (or the polar slope formula)
\item 1 point: slope $\approx0.637$
\end{itemize}
\textbf{(d)} $\frac{dr}{d\theta}=1+2\cos2\theta=0$ at $\theta=\frac\pi3,\frac{2\pi}{3}$. $r(0)=0$, $r(\frac\pi3)\approx1.913$, $r(\frac{2\pi}3)\approx1.228$, $r(\pi)=\pi$; all are nonnegative, so distance $=r$. The greatest distance is $\pi\approx3.142$.
\begin{itemize}[leftmargin=1.6em,itemsep=0pt]
\item 1 point: critical points from $\frac{dr}{d\theta}=0$: $\theta=\frac\pi3$ and $\frac{2\pi}3$
\item 1 point: values of $r$ at the critical points and endpoints
\item 1 point: answer $\pi$ (at $\theta=\pi$)
\end{itemize}
\subsection*{Question 3: A car's velocity (9 points)}
\textbf{(a)} $\frac{10+16}{2}(2)+\frac{16+22}{2}(3)+\frac{22+18}{2}(3)+\frac{18+6}{2}(4)=26+57+60+48=191$. It approximates the distance in meters the car travels during $0\le t\le12$ seconds (the car moves forward, so displacement is distance).
\begin{itemize}[leftmargin=1.6em,itemsep=0pt]
\item 1 point: trapezoidal sum $191$
\item 1 point: meaning: the distance in meters the car travels from $t=0$ to $t=12$ (velocity is positive)
\end{itemize}
\textbf{(b)} $v$ is differentiable, hence continuous on $[2,5]$, and $16<20<22$, so by the IVT there is a $t$ in $(2,5)$ with $v(t)=20$.
\begin{itemize}[leftmargin=1.6em,itemsep=0pt]
\item 1 point: $v$ is continuous because it is differentiable
\item 1 point: $v(2)=16<20<22=v(5)$, so yes by the IVT
\end{itemize}
\textbf{(c)} $\frac{6-18}{4}=-3$. $v$ is differentiable (so continuous), so by the MVT there is a $t$ in $(8,12)$ with $v'(t)=a(t)=-3$.
\begin{itemize}[leftmargin=1.6em,itemsep=0pt]
\item 1 point: $\frac{v(12)-v(8)}{12-8}=-3$
\item 1 point: yes by the MVT, citing differentiability
\end{itemize}
\textbf{(d)} $v(t)=v(12)+\int_{12}^t(-2)\,ds=6-2(t-12)$, so $v(15)=0$. $v\ge0$ on $[12,15]$, so the distance is $\int_{12}^{15}(6-2(t-12))\,dt=18-9=9$ meters.
\begin{itemize}[leftmargin=1.6em,itemsep=0pt]
\item 1 point: $v(t)=6-2(t-12)$, so $v(15)=0$
\item 1 point: distance $\int_{12}^{15}v(t)\,dt=9$ meters
\end{itemize}
\textbf{(e)} $\frac{1}{12}\int_0^{12}v(t)\,dt\approx\frac{191}{12}\approx15.917$ meters per second.
\begin{itemize}[leftmargin=1.6em,itemsep=0pt]
\item 1 point: average velocity $\approx\frac{191}{12}\approx15.917$ m/s
\end{itemize}
\subsection*{Question 4: An accumulation function (9 points)}
\textbf{(a)} $g(-3)=5-\int_{-3}^0f=5-(-1+0)=6$ (triangle of area 1 below the axis on $[-3,-2]$; on $[-2,0]$ the areas $-1$ and $+1$ cancel). $g(4)=5+\int_0^4f=5+4+0=9$ (rectangle of area 4; on $[2,4]$ the areas $+1$ and $-1$ cancel).
\begin{itemize}[leftmargin=1.6em,itemsep=0pt]
\item 1 point: $g(-3)=6$
\item 1 point: $g(4)=9$
\end{itemize}
\textbf{(b)} $g'=f$. $f$ changes from negative to positive at $x=-1$: relative minimum. $f$ changes from positive to negative at $x=3$: relative maximum.
\begin{itemize}[leftmargin=1.6em,itemsep=0pt]
\item 1 point: relative minimum at $x=-1$: $g'=f$ changes from negative to positive
\item 1 point: relative maximum at $x=3$: $g'=f$ changes from positive to negative
\end{itemize}
\textbf{(c)} Candidates: endpoints and critical points $x=-1,3$. $g(-3)=6$, $g(-1)=5-\int_{-1}^0f=5-1=4$, $g(3)=5+4+1=10$, $g(4)=9$. The absolute minimum value is $4$.
\begin{itemize}[leftmargin=1.6em,itemsep=0pt]
\item 1 point: candidates: $x=-3,-1,3,4$
\item 1 point: values $6$, $4$, $10$, $9$
\item 1 point: absolute minimum $4$ at $x=-1$
\end{itemize}
\textbf{(d)} $g(3)=10$, so the form is $\frac00$. L'Hôpital: $\lim\frac{f(x)}{2(x-3)}$, again $\frac00$ since $f(3)=0$. Again: $\lim\frac{f'(x)}{2}=\frac{-2}{2}=-1$ ($f$ has slope $-2$ on $(2,4)$).
\begin{itemize}[leftmargin=1.6em,itemsep=0pt]
\item 1 point: shows $\frac00$ ($g(3)=10$) and applies L'Hôpital's rule
\item 1 point: limit $-1$
\end{itemize}
\subsection*{Question 5: A differential equation (9 points)}
\textbf{(a)} At $(0,-1)$: $\frac{dy}{dx}=(1)(1)=1$, so $y=-1+x$ and $f(0.2)\approx-0.8$.
\begin{itemize}[leftmargin=1.6em,itemsep=0pt]
\item 1 point: tangent line $y=-1+x$
\item 1 point: approximation $f(0.2)\approx-0.8$
\end{itemize}
\textbf{(b)} Step 1: slope $(-1)^2(1)=1$, so $f(0.1)\approx-1+0.1(1)=-0.9$. Step 2: slope $(-0.9)^2(1.2)=0.972$, so $f(0.2)\approx-0.9+0.1(0.972)=-0.8028$.
\begin{itemize}[leftmargin=1.6em,itemsep=0pt]
\item 1 point: first step: $f(0.1)\approx-0.9$
\item 1 point: second step: $f(0.2)\approx-0.8028$
\end{itemize}
\textbf{(c)} $\frac{dy}{y^2}=(2x+1)\,dx$, so $-\frac1y=x^2+x+C$. At $(0,-1)$: $1=C$. $-\frac1y=x^2+x+1$, so $y=-\frac{1}{x^2+x+1}$.
\begin{itemize}[leftmargin=1.6em,itemsep=0pt]
\item 1 point: separates variables
\item 1 point: antiderivative $-\frac1y$
\item 1 point: antiderivative $x^2+x$
\item 1 point: uses the initial condition: $C=1$
\item 1 point: solves: $y=-\frac{1}{x^2+x+1}$
\end{itemize}
\subsection*{Question 6: A Taylor polynomial (9 points)}
\textbf{(a)} $P_3(x)=1-2x+\frac{6}{2!}x^2-\frac{24}{3!}x^3=1-2x+3x^2-4x^3$.
\begin{itemize}[leftmargin=1.6em,itemsep=0pt]
\item 1 point: first two terms $1-2x$
\item 1 point: terms $3x^2-4x^3$
\end{itemize}
\textbf{(b)} $P_3(0.5)=1-1+0.75-0.5=0.25$. $|f(0.5)-P_3(0.5)|\le\frac{60}{4!}(0.5)^4=\frac{60}{24}\cdot\frac{1}{16}=0.15625<0.16$.
\begin{itemize}[leftmargin=1.6em,itemsep=0pt]
\item 1 point: $P_3(0.5)=0.25$
\item 1 point: Lagrange form $\frac{M}{4!}(0.5)^4$ with $M=60$
\item 1 point: bound $0.15625<0.16$
\end{itemize}
\textbf{(c)} $\int_0^x(1-2t+3t^2-4t^3)\,dt=x-x^2+x^3-x^4$.
\begin{itemize}[leftmargin=1.6em,itemsep=0pt]
\item 1 point: integrates term by term
\item 1 point: $x-x^2+x^3-x^4$
\end{itemize}
\textbf{(d)} Substitute $x^2$ into $1-2x+3x^2-\cdots$: $1-2x^2+3x^4-\cdots$.
\begin{itemize}[leftmargin=1.6em,itemsep=0pt]
\item 1 point: substitutes $x^2$ for $x$ in the series for $f$
\item 1 point: $1-2x^2+3x^4$
\end{itemize}
\end{document}